% ─────────────────────────────────────────────────────────────────────────────
%  introduccion.tex  —  Introducción general del documento
% ─────────────────────────────────────────────────────────────────────────────

\chapter*{Introducción}
\addcontentsline{toc}{chapter}{Introducción}  % aparece en TOC aunque sea sin número
\markboth{Introducción}{Introducción}

% Cambiar numeración a arábiga a partir de aquí
\pagenumbering{arabic}
\setcounter{page}{1}

El desarrollo de software moderno se enfrenta continuamente a desafíos de diseño que,
aunque varían en dominio y escala, comparten una naturaleza estructural recurrente.
Para responder a estos desafíos de manera sistemática y reutilizable, la comunidad de
ingeniería de software adoptó el concepto de \textbf{patrones de diseño}, popularizado
por Gamma, Helm, Johnson y Vlissides en su obra clásica \textit{Design Patterns: Elements
of Reusable Object-Oriented Software} (1994) \cite{GoF1994}, conocida coloquialmente como
\textit{«el libro de la Banda de los Cuatro» (GoF)}.

Un patrón de diseño es una \emph{solución general y reutilizable} a un problema
recurrente dentro de un contexto determinado. No se trata de código listo para
copiar y pegar, sino de una plantilla o esquema conceptual que debe adaptarse a
cada situación concreta.

\section*{Contexto de aplicación: Airbnb}

A lo largo de este documento, los 23~patrones del catálogo GoF se analizan y aplican
al ecosistema de \textbf{Airbnb}, una de las plataformas de alojamiento entre pares
(\textit{peer-to-peer}) más grandes del mundo. Airbnb conecta a \emph{anfitriones}
que ofrecen espacios habitacionales con \emph{huéspedes} que buscan alojamiento
temporal, gestionando además pagos, reseñas, búsquedas geoespaciales, notificaciones
y políticas de cancelación, entre otros servicios.

Esta plataforma resulta un escenario idóneo para el estudio de patrones de diseño
porque presenta problemáticas reales y complejas:

\begin{itemize}[leftmargin=1.5cm]
  \item \textbf{Creación dinámica de objetos:} distintos tipos de alojamiento
        (apartamento, cabaña, habitación privada, etc.) con atributos y reglas
        de negocio propios.
  \item \textbf{Integración con servicios externos:} pasarelas de pago, proveedores
        de mapas, sistemas de correo y notificaciones push.
  \item \textbf{Gestión del estado:} reservas que transitan por estados como
        \textit{pendiente}, \textit{confirmada}, \textit{activa} y \textit{finalizada}.
  \item \textbf{Escalabilidad y rendimiento:} millones de listados simultáneos que
        exigen estructuras de datos eficientes y acceso controlado a recursos.
  \item \textbf{Comportamiento flexible:} algoritmos de búsqueda, filtrado y
        recomendación que varían según preferencias del usuario.
\end{itemize}

\section*{Organización del documento}

Los patrones se presentan agrupados en las tres categorías canónicas del GoF:

\begin{enumerate}[leftmargin=1.5cm]
  \item \textbf{Patrones Creacionales} (5 patrones): abordan la forma en que se
        crean los objetos, desacoplando el código cliente de las clases concretas
        que instancia.
  \item \textbf{Patrones Estructurales} (7 patrones): describen cómo se componen
        clases y objetos para formar estructuras más grandes y flexibles.
  \item \textbf{Patrones de Comportamiento} (10 patrones): se ocupan de los
        algoritmos y la asignación de responsabilidades entre objetos.
\end{enumerate}

Dentro de cada categoría los patrones se presentan en orden alfabético. Para cada uno
se incluye: la situación-problema en el contexto de Airbnb, el diagrama de clases UML,
la descripción de roles, un \textit{main} de ejemplo con su salida esperada, el código
fuente de las clases y las conclusiones particulares del patrón.

\bigskip